\documentclass[11pt]{article}
\usepackage[margin=1in]{geometry}
\usepackage{amsmath,amssymb,amsthm,booktabs,hyperref}
\hypersetup{colorlinks=true,urlcolor=blue,linkcolor=blue}
\setlength{\emergencystretch}{1em}
\newtheorem{theorem}{Theorem}
\newcommand{\Phiuni}{\Phi^{\mathrm{uni}}}
\title{A unimodal cyclotomic generating function\protect\\with no prime-index cyclotomic factor}
\author{Principia Math In Progress}
\date{27 September 2026}
\begin{document}
\maketitle

\paragraph{Task 1 (MathDB 358386).}
Billey and Swanson's Conjecture 48 in \cite{BS} asks whether every
nonconstant unimodal basic cyclotomic generating function has a factor
$\Phi_p(q)$ for some prime $p$. A basic cyclotomic generating function is a
monic polynomial with constant term 1 and nonnegative integer coefficients
whose roots are roots of unity. Their monoid of unimodal such polynomials
is denoted by $\Phiuni$. We disprove the conjecture with an explicit
degree-66 polynomial. The statement occurs on page 32 of the final
article and page 33 of arXiv:2305.07620v3.

\begin{theorem}
The polynomial
\[
 F(q)=\Phi_4(q)^7\Phi_9(q)^4\Phi_{25}(q)\Phi_{30}(q)
\]
belongs to $\Phiuni$, is nonconstant, and has no factor $\Phi_p(q)$
with $p$ prime.
\end{theorem}

\begin{proof}
The four cyclotomic polynomials in the displayed product are
\begin{align*}
 \Phi_4(q)&=1+q^2,& \Phi_9(q)&=1+q^3+q^6,\\
 \Phi_{25}(q)&=1+q^5+q^{10}+q^{15}+q^{20},\\
 \Phi_{30}(q)&=1+q-q^3-q^4-q^5+q^7+q^8.
\end{align*}
These identities follow, for example, by exact division in
$q^n-1=\prod_{d\mid n}\Phi_d(q)$.
Each factor is reciprocal, so $F$ is reciprocal of degree
$7\cdot2+4\cdot6+20+8=66$. Write $F(q)=\sum_{s=0}^{66}c_sq^s$.
The coefficients through the midpoint are as follows; the others are
given by $c_{66-s}=c_s$.
\begin{center}
\begin{tabular}{rr@{\qquad}rr@{\qquad}rr@{\qquad}rr}
\toprule
$s$&$c_s$&$s$&$c_s$&$s$&$c_s$&$s$&$c_s$\\
\midrule
0&1&9&198&18&920&27&1450\\
1&1&10&283&19&940&28&1570\\
2&7&11&337&20&1025&29&1588\\
3&10&12&437&21&1058&30&1704\\
4&24&13&524&22&1109&31&1737\\
5&42&14&592&23&1184&32&1746\\
6&63&15&701&24&1220&33&1810\\
7&105&16&763&25&1319&&\\
8&147&17&829&26&1381&&\\
\bottomrule
\end{tabular}
\end{center}
\newpage
For completeness, the following finite sum specifies every entry of the
table directly. Define
\[
b_s=\sum_{i=0}^{7}\sum_{j=0}^{4}\sum_{k=0}^{4-j}\sum_{h=0}^{4}
 {7\choose i}\frac{4!}{j!k!(4-j-k)!}
 \mathbf{1}_{\{s=2i+3j+6k+5h\}},
\]
so $b_s$ is the coefficient of $q^s$ in
$(1+q^2)^7(1+q^3+q^6)^4(1+q^5+q^{10}+q^{15}+q^{20})$,
with $b_s=0$ outside $0\le s\le58$. Multiplication by $\Phi_{30}$ gives
\[
c_s=b_s+b_{s-1}-b_{s-3}-b_{s-4}-b_{s-5}+b_{s-7}+b_{s-8}.
\]
This proves the tabulated values by finite integer arithmetic. They satisfy
$1=c_0\le c_1\le\cdots\le c_{33}$. Reciprocity now proves that every
coefficient is positive and that the entire sequence is unimodal. The
factorization also proves that $F$ is monic with constant term 1 and that
every root is a root of unity. Thus $F\in\Phiuni$.

Finally, every root of $F$ has exact order in $\{4,9,25,30\}$.
If $\Phi_p$ divided $F$ for a prime $p$, a primitive $p$th root of unity
would be a root of $F$. None of the four possible orders is prime,
which is a contradiction. Since $\deg F=66>0$, this is a counterexample
to the stated conjecture.
\end{proof}

\paragraph{Verification and scope.}
The accompanying standard-library Python verifier constructs cyclotomic
polynomials from the divisor identity, checks the table by two exact
expansions, and checks nondivisibility by $\Phi_p$ for every prime
$p\le67$ by polynomial division. Larger primes are excluded by degree.
An additional SymPy expansion independently reproduced all 67 coefficients.
No claim is made that degree 66 is minimal. This example is not log-concave:
$c_1^2=1<7=c_0c_2$. The source's prime-power factor theorem is consistent
with the example, which has factors $\Phi_4,\Phi_9,\Phi_{25}$.

\paragraph{Proposed Tasks 2--5.}
The following directions are proposed after the completed resolution of
Task 1; they are not claimed to be solved here, and their novelty has not
been audited.
\begin{enumerate}
\setcounter{enumi}{1}
\item Determine the smallest degree of a nonconstant member of $\Phiuni$
having no prime-index cyclotomic factor, with a certificate excluding
every smaller degree.
\item Characterize the nonnegative integer quadruples $(a,b,c,d)$ for
which $\Phi_4^a\Phi_9^b\Phi_{25}^c\Phi_{30}^d$ is unimodal.
\item Give an algorithm that, for an input degree bound $D$, enumerates
exactly the irreducible elements of $\Phiuni$ of degree at most $D$
that have no prime-index cyclotomic factor, and prove its correctness.
\item Let $N(D)$ count the nonconstant elements of $\Phiuni$ of degree
at most $D$ having no prime-index cyclotomic factor. Determine the
asymptotic order of $\log N(D)$ as $D\to\infty$.
\end{enumerate}

\begin{thebibliography}{1}
\bibitem{BS} S. C. Billey and J. P. Swanson,
\emph{Cyclotomic Generating Functions}, Electronic Journal of
Combinatorics \textbf{31}(4) (2024), P4.4, 43 pp.
\href{https://doi.org/10.37236/12687}{doi:10.37236/12687}.
Current manuscript: \href{https://arxiv.org/abs/2305.07620}{arXiv:2305.07620v3}.
\end{thebibliography}
\end{document}
